\clearpage
\section{Hadronic Event Selection\label{sec:eventSelection}}
As discused in [REF], all-hadronic SUSY signatures consist of events with
no isolated and detectable objects apart from energetic jets. To search
for an excess in such events, it is convenient to use a jet-based variable,
\scalht, to quantify the energy in an event.  \scalht is defined as the scalar
sum of the transverse energy of the jets in the event. The challanges due to large
backgrounds in a search for an excess in all-hadronic events becomes evident 
when compairing observed data events overlaid with simulated events from 
SM processes as a function of \scalht Figure [REF].
 
The dominant background are azimuthally-balanced multi-jet events, stemming from QCD
processes. This search reduces this background to negligible levels in the
selected region by employing the \alphat variable further discussd in section [REF].
Furthermore, by construction, \alphat also reduces backgrounds from severely mis-measured
jets. 

In absence of the multi-jet background, the remaining significant backgrounds are ...

\subsection{Data}
The data analyzed was recorded by CMS in 2012 between April 5th and Dec. 5th and
totals 19.47$\pm$ 2.6\%. The collected data is certified on a run-by-run 
basis, where initial automatic certification requires the LHC beams to be declared
stable and all CMS subdetectors ON. Further monitoring of the data was done realtime
by experts of each subdetector trough analysis of histograms updated and filled
each lumi section. Final cerfitication was done offline, and lumi sections passing
all criteria were listed in a Golden JSON file to be used by all analyses. 

\subsection{Event quality}

Each event is subjected to a series of commonly used filters in CMS to ensure good
quality data. Minimal requirements are that at least one primary vertex is identified 
and 25\% of the reconstructed tracks to be of good quality. Addionally, various filters 
prescribed by the MET group[REF] are applied. Events containing muons with inconsistent
energy are flagged by the muon pog [REF] and filtered out of the analysis. 

\subsection{Triggers}

\subsubsection{Hadronic signal region and control samples\label{sec:signal_triggers}} 

Only events passing one or more HLT triggers based on online quantanties 
are recorded to be analyzed. For any analysis, in general, it is 
not expected that all recorded events reconstructed offline, pass the online
trigger as detector conditions, energy corrections, and object-based quantaties
differ offline. In this analysis, cross-triggers at the HLT
based on quantities \scalht and \alphat (labelled as \verb!HTxxx_AlphaT0pyy!) 
are used with various thresholds to record candidate events for the hadronic (signal)
region. In order to keep the cross-trigger's computational time low, the online quantaties
are constructed using simple calorimeter based jets (calo jets), and the use of
partcle-flow jets in this analysis is expected to introduce inefficiencies.

Each \scalht bin is seeded by a single trigger chosen based on the
efficiency of the trigger in that \scalht bin. The \alphat thresholds of the
\verb!HTxxx_AlphaT0pyy! triggers were tuned according to the threshold
on the \scalht leg in order to fully suppress QCD multijet events
and simultaneously satisfying other criteria, such as maintaining
acceptable trigger rates.


%Table~\ref{tab:htalphat-triggers} summarises the thresholds used for
%the \verb!HTxxx_AlphaT0pyy! triggers. 
%
The \verb!HTxxx_AlphaT0pyy! trigger efficiencies are measured with a
reference (\ie, unbiased) event sample recorded by an unprescaled,
loosely-isolated, eta-restricted single muon trigger,
\verb!HLT_IsoMu24_eta2p1!, within the \verb!SingleMu! dataset. A
sample of events containing at least one isolated muon with $\pt >
25\gev$ and $|\eta| < 2.1$ is used (similar to the \mj control sample
defined in Section~\ref{sec:def-control-samples}). A cut of $\Delta
{\rm R} > 0.5$ is placed between all muons and jets in each event, and
only jets are considered in the calculation of \scalht, \mht, and
\alphat, \ie the muon is ignored.
%
%Table~\ref{tab:htalphat-triggers} summarises the measured efficiencies
%for the \verb!HTxxx_AlphaT0pyy! triggers in the relevant \scalht
%bins. The trigger efficiencies are measured for both \njet
%multiplicity bins. The efficiencies are generally observed to be high,
%$\sim$100\%, except for the low \scalht region $200 < \scalht <
%325\gev$. The inefficiencies at low \scalht are mainly due to the L1
%seeds for which thresholds were raised to a relatively high level in
%order to maintain trigger rates in the high-pileup conditions towards
%the end of Run 1. The inefficiencies are slightly larger in the higher
%jet multiplicity category due to a larger number of jets summing to
%the same \scalht, resulting in softer
%jets. Figures~\ref{fig:eff-alphat-le3j} and~\ref{fig:eff-alphat-ge4j}
%show the efficiency curves for the \verb!HTxxx_AlphaT0pyy! triggers in
%the three lowest \scalht bins, for the \njetlow and \njethigh
%categories, respectively. The efficiencies are also determined at the
%level of event categories and no dependence on \nb is observed, with
%efficiencies agreeing within statistical uncertainties.
%

\begin{figure}[!h]
  \begin{center}
    \subfigure[Differential, $375 < \scalht < 475 \gev$]{
      \includegraphics[width=0.4\textwidth,page=20]{figures/trigger/plotDump/v29/HT375_475_100_100_50_AlphaT_HT300xaT0p53_PF_le3j_RunAtFNAL.pdf}
    }
    \subfigure[Cumulative, $375 < \scalht < 475 \gev$]{
      \includegraphics[width=0.4\textwidth,page=30]{figures/trigger/plotDump/v29/HT375_475_100_100_50_AlphaT_HT300xaT0p53_PF_le3j_RunAtFNAL.pdf}
    } \\
    \subfigure[Differential, $475 < \scalht < 525 \gev$]{
      \includegraphics[width=0.4\textwidth,page=20]{figures/trigger/plotDump/v29/HT475_575_100_100_50_AlphaT_HT350xaT0p52_PF_le3j_RunAtFNAL.pdf}
    } 
    \subfigure[Cumulative, $475 < \scalht < 525 \gev$]{
      \includegraphics[width=0.4\textwidth,page=30]{figures/trigger/plotDump/v29/HT475_575_100_100_50_AlphaT_HT350xaT0p52_PF_le3j_RunAtFNAL.pdf}
    } \\
    \subfigure[Differential, $525 < \scalht < 675 \gev$]{
      \includegraphics[width=0.4\textwidth,page=18]{figures/trigger/plotDump/v29/HT575_675_100_100_50_AlphaT_HT400xaT0p51_PF_le3j_RunAtFNAL.pdf}
    }
    \subfigure[Cumulative, $525 < \scalht < 675 \gev$]{
      \includegraphics[width=0.4\textwidth,page=27]{figures/trigger/plotDump/v29/HT575_675_100_100_50_AlphaT_HT400xaT0p51_PF_le3j_RunAtFNAL.pdf}
    } \\
    \caption{\label{fig:eff-alphat-le3j}
      (Left) Differential and (Right) cumulative efficiency turn-on 
      curves for the \scalht-\alphat cross triggers (as summarised in 
      Table~\ref{tab:htalphat-triggers}) that record events for the
      three lowest \scalht bins  for events satisfying \njetlow. 
    }
  \end{center}
\end{figure}
%
\begin{figure}[!h]
  \begin{center}
    \subfigure[Differential, $375 < \scalht < 475 \gev$]{
      \includegraphics[width=0.4\textwidth,page=20]{figures/trigger/plotDump/v29/HT375_475_100_100_50_AlphaT_HT300xaT0p53_PF_ge4j_RunAtFNAL.pdf}
    }
    \subfigure[Cumulative, $375 < \scalht < 475 \gev$]{
      \includegraphics[width=0.4\textwidth,page=30]{figures/trigger/plotDump/v29/HT375_475_100_100_50_AlphaT_HT300xaT0p53_PF_ge4j_RunAtFNAL.pdf}
    } \\
    \subfigure[Differential, $475 < \scalht < 525 \gev$]{
      \includegraphics[width=0.4\textwidth,page=20]{figures/trigger/plotDump/v29/HT475_575_100_100_50_AlphaT_HT350xaT0p52_PF_ge4j_RunAtFNAL.pdf}
    } 
    \subfigure[Cumulative, $475 < \scalht < 525 \gev$]{
      \includegraphics[width=0.4\textwidth,page=30]{figures/trigger/plotDump/v29/HT475_575_100_100_50_AlphaT_HT350xaT0p52_PF_ge4j_RunAtFNAL.pdf}
    } \\
    \subfigure[Differential, $525 < \scalht < 675 \gev$]{
      \includegraphics[width=0.4\textwidth,page=18]{figures/trigger/plotDump/v29/HT575_675_100_100_50_AlphaT_HT400xaT0p51_PF_ge4j_RunAtFNAL.pdf}
    }
    \subfigure[Cumulative, $525 < \scalht < 675 \gev$]{
      \includegraphics[width=0.4\textwidth,page=27]{figures/trigger/plotDump/v29/HT575_675_100_100_50_AlphaT_HT400xaT0p51_PF_ge4j_RunAtFNAL.pdf}
    } \\
    \caption{\label{fig:eff-alphat-ge4j}
      (Left) Differential and (Right) cumulative efficiency turn-on 
      curves for the \scalht-\alphat cross triggers (as summarised in 
      Table~\ref{tab:htalphat-triggers}) that record events for the
      three lowest \scalht bins  for events satisfying \njethigh. 
    }
  \end{center}
\end{figure}


\subsubsection{Muon control samples\label{sec:muon_triggers}}

%The procedure to determine the muon trigger efficiency follows the
%recommendations of the muon POG~\cite{ref:muon-eff}. Events for the
%\mj and \mmj control samples are recorded with a loosely-isolated,
%$\eta$-restricted muon trigger \verb!HLT_IsoMu24_eta2p1!. The muon
%trigger efficiency is determined for each of the \mj and \mmj control
%samples according to the binning scheme in \njet and \scalht. Events
%with muons satisfying the acceptance requirements $\pt > 25\gev$ and
%$|\eta| < 2.1$ are considered.
%
%The \verb!HLT_IsoMu24_eta2p1! trigger efficiency is determined from
%data in bins of muon $\pt$ and $|\eta|$ by the muon POG from a
%tag-and-probe method~\cite{ref:muon-eff}. These efficiencies are
%weighted by MC yields binned according to the muon $\pt$ and $|\eta|$
%in order to determine a single weighted efficiency measurement per
%(\njet,\scalht) bin. Simulation-to-data scale factors for muon
%identification and isolation efficiencies are applied to MC
%events. These scale factors are also determined by the muon POG in
%bins of muon $\pt$ and $|\eta|$. Table~\ref{tab:muon-effs} summarises
%the muon trigger efficiencies which are assumed to have a relative
%systematic uncertainty of 1\%~\cite{ref:muon-eff}. Efficiencies are
%higher for the \mmj sample due to the fact that both muons must
%satisfy $\pt > 30\gev$ and so can be the source of the positive
%trigger decision. A further consequence of these thresholds
%requirements is that there is little dependence on \scalht. 
%
%\begin{table}[!h]
%  \caption{Muon trigger efficiencies (\%). Statistical uncertainties
%    are at the per-mille level, while a relative systematic
%    uncertainty on all measurements is assumed to be 1\%.}  
%  \label{tab:muon-effs}
%  \centering
%  \footnotesize
%  \begin{tabular}{ cccccc }
%    \hline
%    \hline
%    \scalht (GeV) \textbackslash \njet & \multicolumn{2}{c}{\mj} & \multicolumn{2}{c}{\mmj} \\ [0.5ex]
%                                       & 2-3                     & $\geq$4 & 2-3 & $\geq$4  \\ [0.5ex]
%    \hline
%%    200--275  & 89.1 & 89.8 & 98.9 & 98.8 \\
%%    275--325  & 89.3 & 89.8 & 98.9 & 98.8 \\
%%    325--375  & 89.5 & 90.0 & 98.9 & 98.8 \\
%%    375--475  & 89.7 & 90.3 & 98.9 & 98.9 \\
%%    475--575  & 89.8 & 90.5 & 98.9 & 98.9 \\
%%    575--675  & 90.0 & 90.6 & 98.9 & 98.9 \\
%%    675--775  & 90.1 & 90.7 & 99.0 & 98.9 \\
%%    775--875  & 90.2 & 90.8 & 98.9 & 98.9 \\
%%    875--975  & 90.4 & 90.6 & 99.0 & 98.9 \\
%%    975--1075 & 90.3 & 90.6 & 99.0 & 98.9 \\
%%    $>$1075   & 90.0 & 91.2 & 99.0 & 99.1 \\
%    150--200  & 87.2 & 88.1 & 98.4 & 98.4  \\
%    200--275  & 87.5 & 88.1 & 98.5 & 98.4  \\
%    275--325  & 87.8 & 88.2 & 98.5 & 98.4  \\
%    325--375  & 87.9 & 88.4 & 98.6 & 98.6  \\
%    375--475  & 88.1 & 88.6 & 98.6 & 98.5  \\
%    475--575  & 88.2 & 88.8 & 98.6 & 98.6  \\
%    575--675  & 88.4 & 88.9 & 98.6 & 98.6  \\
%    675--775  & 88.5 & 89.0 & 98.7 & 98.6  \\
%    775--875  & 88.6 & 89.1 & 98.6 & 98.6  \\
%    875--975  & 88.8 & 89.0 & 98.7 & 98.6  \\
%    975--1075 & 88.7 & 89.0 & 98.7 & 98.8  \\
%    $>$1075   & 88.4 & 89.6 & 98.7 & 98.7  \\
%    \hline
%    \hline
%  \end{tabular}
%\end{table}
%
%\subsection{Photon control sample\label{sec:photon_triggers}}
%
%Events for the photon control sample are recorded with the
%\verb!HLT_Photon150! trigger, which is $\sim100\%$ efficient for
%$E_{\rm T}^{\rm photon} > 165\gev$ and $\scalht > 375\gev$, as shown
%in Figure~\ref{fig:eff-photon}. The efficiency measurement is made
%using the \verb!HLT_Photon90! trigger as a reference, and we assume
%that the \verb!L1_SingleEG22! seed is fully efficient.
%
%\begin{figure}[!h]
%  \begin{center}
%    \subfigure[\njetlow]{
%      \includegraphics[width=0.48\textwidth,page=3,trim=40 50 160 120,clip=true]{figures/trigger/g_barrel_375_caloJet_le3j.pdf}
%    }
%    \subfigure[\njethigh]{
%      \includegraphics[width=0.48\textwidth,page=3,trim=40 50 160 120,clip=true]{figures/trigger/g_barrel_375_caloJet_ge4j.pdf}
%    } \\
%    \caption{\label{fig:eff-photon} Efficiency turn-on curves for
%      the \texttt{HLT\_Photon150} trigger that records events that
%      satisfy the \gj selection criteria, $E_{\rm T}^{\rm photon} >
%      165\gev$, $\scalht > 375\gev$, and \njetlow (Left) and \njethigh
%      (Right). }
%  \end{center}
%\end{figure}

\subsection{event selection}

%Once all the signal selection requirements have been imposed, the
%contribution from QCD multijet events is expected to be negligible, as
%demonstrated in Section~\ref{sec:qcd}.
%
%The remaining backgrounds in the hadronic signal region stem from SM
%processes with genuine \met in the final state.  In the case of events
%where no b-quark jets are identified, the largest backgrounds with
%genuine \met arise from the production of W and Z bosons in
%association with jets. The weak decay \znunu\ is the only significant
%contribution from Z + jets events. For W + jets events, the two
%dominant sources are leptonic W decays in which the lepton is not
%reconstructed or fails the isolation or acceptance requirements, and
%the weak decay $\wtaunu$ where the $\tau$ decays hadronically and is
%identified as a jet. Contributions from SM processes such as
%single-top, Drell-Yan, and diboson production are also expected. For
%events with one or more reconstructed b-quark jets, \ttbar production
%followed by semi-leptonic weak decays becomes the most important
%single background source. For events with only one reconstructed
%b-quark jet, the contribution of both W + jets and Z + jets
%backgrounds are of a similar size to the \ttbar background.  For
%events with two reconstructed b-quark jets, \ttbar production
%dominates, while events with three or more reconstructed b-quark jets
%originate almost exclusively from \ttbar events, in which at least one
%jet is misidentified as originating from a bottom quark.
%
%In order to estimate the contributions from each of these backgrounds,
%three data control samples are used, which are binned identically to
%the signal region. Two independent estimates of the irreducible
%background of \znunu\ + jets events are determined from the data
%control samples comprising $Z\rightarrow\mu\mu$ + jets and $\gamma$ +
%jets events, which have similar kinematic properties but different
%acceptances. The $Z\rightarrow\mu\mu$ + jets events have similar
%kinematic properties when the two muons are ignored, but a smaller
%branching fraction, while the $\gamma$ + jets events have similar
%kinematic properties when the photon is
%ignored~\cite{PAS-SUS-08-002,Bern:2011pa}, but a larger production
%cross section. A $\mu$ + jets data sample provides an estimate for all
%other SM backgrounds, which is dominated by \ttbar and W production
%leading to W + jets final states.
%
%As described previously, the event selection criteria for the control
%samples are defined to ensure that any potential contamination from
%multijet events is negligible. Further, the control sample selection
%criteria also suppress contributions from a wide variety of SUSY
%models, including those considered in this analysis. Any potential
%signal contamination in the data control samples is accounted for in
%the fitting procedure described in Section~\ref{sec:results}.
%
%\subsection{Overview of the method\label{sec:background-method}}
%
%The method used to estimate the aforementioned SM background
%contributions in the hadronic signal region relies on the use of a
%{\it transfer factor} (TF) determined from MC samples to transform the
%observed yield in a given \scalht, jet (\njet) and b-tag (\nb)
%multiplicity bin of a control sample, $\nobs^{\rm
%  control}(\scalht,\njet,\nb)$, into a predicted yield for the
%corresponding bin of the hadronic signal region, $\npre^{\rm
%  signal}(\scalht,\njet,\nb)$. The choice of \njet and \nb event
%categorisation and \scalht binning in the control samples is identical
%to that for the signal region, as defined in Table~\ref{tab:ht-bins}
%in Section~\ref{sec:selection}. 
%
%Each transfer factor is simply a ratio of the yields obtained from MC
%simulation for the same bin of the signal region and a given control
%sample:
%
%\begin{equation}
%  \label{equ:tf-ratio}
%  {\rm TF} = \frac{N_{\rm MC}^{\rm signal}(\scalht,\njet,\nb)}{N_{\rm
%      MC}^{\rm control}(\scalht,\njet,\nb)} 
%\end{equation}
%
%In this way, ``na\"ive'' predictions for the total SM background can
%be made by considering separately the sum of the predictions from the
%\mj and \gj samples or the \mj and \mmj samples:
%
%\begin{equation}
%  \label{equ:pred-method}
%  \npre^{\rm signal}(\scalht,\njet,\nb) = \frac{N_{\rm MC}^{\rm
%      signal}(\scalht,\njet,\nb)}{N_{\rm MC}^{\rm
%      control}(\scalht,\njet,\nb)} \times \nobs^{\rm
%    control}(\scalht,\njet,\nb)   
%\end{equation}
%
%When constructing the transfer factors, the MC expectations for the
%following SM processes are considered: W + jets ($N_{\rm W}$), \ttbar
%+ jets ($N_{\ttbar}$), \znunu\ + jets ($N_{\znunu}$), DY + jets
%($N_{\mathrm DY}$), \gj ($N_\gamma$), single top + jets
%production via the s, t, and tW-channels ($N_{\rm top}$), and WW +
%jets, WZ + jets, and ZZ + jets ($N_{\rm di-boson}$). Details on the MC
%samples used are given in Sec.~\ref{sec:mc-samples}. All MC samples
%are normalised to the intergrated luminosity of the appopriate data
%sample.
%
%While ``na\"ive'' predictions can be made using
%Equ.~\ref{equ:pred-method}, the fitting prcedure that provides the
%final result is defined formally by the likelihood model described in
%Sec.~\ref{sec:statistics}. In summary, the observation in each bin
%(defined in terms of the variables \njet, \nb, and \scalht) of the
%signal sample is modelled as Poisson-distributed about the sum of a SM
%expectation (and a potential signal contribution). The components of
%this SM expectation are related to the expected yields in the control
%samples via transfer factors derived from simulation, as described
%in Sec.~\ref{sec:background-method}. The observations in each bin
%(again defined by \njet, \nb, and \scalht) of the control samples are
%similarly modelled as Poisson-distributed about the expectated yields
%for each control sample. In this way, for a given bin, the observed
%yields in the signal and control samples are all connected via the
%transfer factors derived from simulation.
%
%The sum of expected yields from all MC samples, obtained for the
%relevant control sample selection, enter the denominator of each
%transfer factor:
%
%\begin{equation}
%  \label{equ:ratio-denom}
%  N_{\rm MC}^{\rm control}(\scalht,\njet,\nb) = N_{\rm W} + N_{\ttbar} + N_{\znunu} +
%N_{\rm DY} + N_{\gamma} + N_{\rm top} + N_{\rm di-boson}
%\end{equation}
%
%For the b jet multiplicity bins satisfying $n_b \leq 1$, the \mj
%control sample is used to predict primarily the W + jets and \ttbar +
%jets backgrounds, but all remaining residual SM backgrounds except
%\znunu\ + jets events are also considered. (The \znunu\ + jets
%background is instead accounted for through the \mmj and \gj samples.)
%The sum of expected yields from all MC samples except \znunu\ + jets,
%obtained for the hadronic signal region selection, enter the numerator
%of each transfer factor:
%
%\begin{equation}
%  \label{equ:ratio-numer-mj}
%  N_{\rm MC}^{\rm signal}(\scalht,\njet,\nb \leq 1) = N_{\rm W} +
%  N_{\ttbar} + N_{\rm DY} + N_{\rm top} + N_{\rm di-boson}
%\end{equation}
%
%For the same b jet multiplicity bins ($0 \leq \nb \leq 1$), the \mmj
%and \gj control samples are used to predict the \znunu\ + jets process
%only, and the expected yields in the bins of the signal region as
%obtained from the \znunu\ sample enter the numerator of each transfer
%factor:
%
%\begin{equation}
%  \label{equ:ratio-numer-mmj}
%  N_{\rm MC}^{\rm signal}(\scalht,\njet,0 \leq \nb \leq 1) = N_{\znunu}
%\end{equation}
%
%For the b jet multiplicity bins satisfying $n_b \geq 2$, the \mj
%control sample is again used to predict primarily the W + jets and
%\ttbar + jets backgrounds, but also to predict all other remaining
%residual SM backgrounds {\it including} \znunu\ + jets events. The sum
%of expected yields from all MC samples {\it including} \znunu\ + jets,
%obtained for the hadronic signal region selection, enter the numerator
%of each transfer factor:
%
%\begin{equation}
%  \label{equ:ratio-numer-mj}
%  N_{\rm MC}^{\rm signal}(\scalht,\njet,\nb \geq 2) = N_{\rm W} +
%  N_{\ttbar} + N_{\rm DY} + N_{\rm top} + N_{\rm di-boson} + N_{\znunu}
%\end{equation}
%
%In this case, the \mmj and \gj control samples are not used, as the
%yields in these two data control samples are expected to be negligible
%due to the requirement of at least two b jets per event. The method of
%using a W + jets sample to predict the \znunu\ + jets background has
%been used previously~\cite{RA1Paper, RA1Paper2011}, and this approach
%is addressed by a set of closure tests described in
%Sec.~\ref{sec:bkgd-syst}, in which a \mj sample (rich in W + jets and
%\ttbar events) is used to make predictions of yields in a \mmj sample
%(rich in Z$\rightarrow\mu\mu$ + jets events). The data control samples
%used to predict the SM backgrounds for each event category are
%summarised in Table~\ref{tab:fit-plots}.
%
%\begin{table}[ht!]
%  \caption{Summary of control samples used to predict the SM
%    background for each event category. }
%  \label{tab:fit-plots}
%  \centering
%  \begin{tabular}{ lll }
%    \hline
%    \hline
%    \njet   & \nb     & Control samples \\ [1.0ex]
%    \hline
%    2--3    & 0       & \mj, \mmj, \gj  \\
%    2--3    & 1       & \mj, \mmj, \gj  \\
%    2--3    & 2       & \mj             \\
%    $\geq$4 & 0       & \mj, \mmj, \gj  \\
%    $\geq$4 & 1       & \mj, \mmj, \gj  \\
%    $\geq$4 & 2       & \mj             \\
%    $\geq$4 & 3       & \mj             \\
%    $\geq$4 & $\geq4$ & \mj             \\
%    \hline
%    \hline
%  \end{tabular}
%\end{table}
%
%The selection criteria for the three control samples closely resemble
%those for the signal region, differing mainly through the use of a
%muon, di-muon, or photon {\it tag} (that is ignored in the calculation
%of jet-based kinematic variables such as \scalht, \mht, \alphat, \etc)
%and some minimal additional kinematic requirements (\eg invariant or
%transerve mass windows) to obtain W, Z, and \ttbar-enriched event
%samples. The same selection criteria are designed to suppress signal
%contamination in the control samples so that unbiased data-driven
%estimates for the SM backgrounds in the signal region can be
%made. Hence, we refer to these samples as {\it control} samples
%although in the final simultaneous fit, any potential signal
%contamination is properly taken into account.
%
%The control sample definitions and binning scheme are chosen so that
%the reliance on simulation to extrapolate correctly from a control
%region to the signal region is minimised. Many systematic effects are
%expected to cancel largely in the transfer factor. However, a
%systematic uncertainty is assigned to each transfer factor to account
%for theoretical uncertainties and effects such as the mismodelling of
%kinematics (\eg acceptances) and instrumental effects (\eg
%reconstruction efficiencies), as described in
%Sec.~\ref{sec:bkgd-syst}.
%
%%Kinematic cuts are applied to enrich as much as possible the \wj,
%%\ttbar, and \znunu components in the muon and di-muon control
%%samples. The definition of the samples are geared towards efficiency
%%rather than purity (even so, the purities are at the level $>$90\%)
%%and any contamination from "backgrounds" (\eg \ttbar in the case of
%%the \mmj sample) are simply incorporated into the transfer
%%factors. Alternatively, a zero b-jet requirement can be applied to the
%%\mmj sample to obtain a higher purity of Z$\rightarrow\mu\mu$ + jets
%%events (\ie, with reduced contamination from \ttbar), which can then
%%be used to give an expectation for the \znunu + jets background for
%%all b-jet categories in the signal region.
%
%\subsection{Definition of the control samples\label{sec:def-control-samples}}
%
%\subsubsection{The \texorpdfstring{\mj}{muon plus jets} control sample}
%
%Events from the \wj and \ttbar processes are found in the hadronic
%signal sample due to unidentified leptons (either out of acceptance or
%not reconstructed) and hadronic tau decays originating from
%high-p$_{T}$ W bosons. An estimate of these background processes is
%obtained through the use of a \mj sample. The selection criteria for
%this sample are chosen to identify W bosons decaying to a muon and a
%neutrino in the phase-space of the signal. The muon is not considered
%in the calculation of event-level variables such as \scalht, \mht and
%\alphat. All cuts on such jet-based quantities are consistent with
%those applied in the hadronic search region and the same \njet, \nb,
%and \scalht binning is used. The only exception is that no \alphat
%requirement is made, as motivated by the discussion in
%Sec.~\ref{sec:larger}. In order to select events containing W bosons,
%exactly one tight isolated muon within an acceptance of \PT $>$ 30
%\gev and $|\eta| <$ 2.1 is required (due to the trigger), and the
%transverse mass of the W candidate must satisfy $30 < \mt(\mu,\pfmet)
%< 125\gev$ (to suppress QCD multijet and potential signal
%events). Events are vetoed if $\Delta R(\mu,\textrm{jet}_i) < 0.5$
%running over all jets $i$. The single isolated track veto, described
%in Sections~\ref{sec:reconstruction} and~\ref{sec:vetoes}, is also
%applied, which considers all single isolated tracks in the event
%except that associated with the identified, isolated muon. Finally,
%the cleaning cut $\mht/\pfmet$ is also applied, as done in the signal
%region, where the \pfmet is adjusted to account for the transverse
%momentum of the identified, isolated muon.
%
%% Events are vetoed if either of the following conditions are met:
%% $\Delta R(\mu,\textrm{jet}_i) < 0.5$, running over all jets $i$; or
%% a second muon candidate exists that is either loose, non-isolated or
%% outside acceptance and the two muons have an invariant mass that
%% satisfies $m_{Z} - 25 < M_{\mu_1\mu_2} < m_{Z} + 25$ (to suppress
%% $Z\rightarrow\mu\mu$).
%
%\subsubsection{The \texorpdfstring{\mmj}{di-muon plus jets} control sample}
%
%The \znunu\ + jets process forms an irreducible background and can be
%estimated using the \zmumu + jets process, which has similar kinematic
%properties but a different acceptance and a smaller branching ratio. A
%background estimate is obtained through the use of a \mmj sample. Most
%of the selection criteria are identical to those for the \mj sample,
%but the few that differ are tuned to identify Z bosons decaying to two
%muons in the kinematic phase space of the signal region. The muons are
%not considered in the calculation of event-level variables such as
%\scalht, \mht and \alphat. All cuts on such jet-based quantities are
%consistent with those applied in the hadronic search region and the
%same \njet, \nb, and \scalht binning is used. The only exception is
%that no \alphat requirement is made, as motivated by the discussion in
%Sec.~\ref{sec:larger}. In order to select an event sample containing Z
%bosons, exactly two tight isolated muons within an acceptance of $\Pt
%> 30\gev$ and $|\eta| < 2.1$ are required (due to the trigger). The
%invariant mass of the two muons must satisfy $m_{Z} - 25 <
%M_{\mu_1\mu_2} < m_{Z} + 25$. Events are vetoed if $\Delta
%R(\mu_{i},\textrm{jet}_j) < 0.5$ is satisfied, running over all muons
%$i$ and all jets $j$. The single isolated track veto, described in
%Sections~\ref{sec:reconstruction} and~\ref{sec:vetoes}, is also
%applied, considering all single isolated tracks in the event except
%those associated with the two identified, isolated muons. Finally, the
%cleaning cut $\mht/\pfmet$ is also applied, as done in the signal
%region, where the \pfmet is adjusted to account for the transverse
%momenta of the two identified, isolated muons. The \mmj sample can be
%used to make predictions in all the \scalht bins, providing coverage
%at low \scalht where the \gj sample cannot.
%
%\subsubsection{The \texorpdfstring{\gj}{photon plus jets} control sample}
%
%The \znunu\ + jets process can also be estimated using the \gj
%process, which has a larger cross section and kinematic properties
%similar to those of \znunu\ events when the photon is
%ignored~\cite{PAS-SUS-08-002,Bern:2011pa}. The \gj sample is defined
%by requiring exactly one photon satisfying tight isolation criteria
%and within an acceptance of $\pt > 165\gev$ and $|\eta| <
%1.45$. Furthermore, events are vetoed if $\Delta
%R(\gamma,\textrm{jet}_j) < 1.0$ is satisfied, running over all jets
%$j$. As for the muon-based samples, the photon is not considered in
%the calculation of event-level variables such as \scalht, \mht and
%\alphat. All cuts on jet-based quantities are consistent with those
%applied in the hadronic search region, and the same \HT binning is
%used. Given that the photon is ignored, the \gj sample can only be
%used for the region $\scalht > 375\gev$ due to the photon acceptance
%of $\pt > 165\gev$ (enforced by the trigger) and the requirement
%$\alphat > 0.55$.
%
%\subsection{Increasing the acceptance of the muon control samples\label{sec:larger}}
%
%As described in Sec.~\ref{sec:def-control-samples} above, the
%selection criteria of the three control samples are defined such that
%the background composition and event kinematics of the three control
%samples mirror as closely as possible those for the signal
%region. This is done in order to minimise the reliance on the
%simulation to model correctly the backgrounds and event kinematics in
%the control and signal samples.
%
%However, in the case of the \mj and \mmj samples, no requirement is
%made on \alphat in the selection criteria of the samples. This is made
%possible by the remaining kinematic selection criteria, which are
%sufficiently selective to ensure that the muon samples remain rich in
%events from the \wj, \ttbar and \zmumu processes with negligible
%contamination from QCD multijet events. These selection criteria
%include, for example, requiring exactly one or two tight isolated
%muon(s), and imposing acceptance windows on the invariant mass of the
%di-muon sytem or the transverse mass of the muon-\pfmet system, as
%described above. 
%%The absence of QCD multijet events is demonstrated by the control
%%distributions shown in Sec.~\ref{sec:est-control-samples} below. 
%Thus, the acceptance of the two muon control samples can be
%significantly increased, which simultaneously improves their
%predictive power and further reduces the effect of any potential
%signal contamination.  In the case of the \gj sample (used only for
%the region $\scalht > 375\gev$), the requirement $\alphat > 0.55$ is
%still necessary to suppress contamination from QCD multijet events,
%even after the substantial photon \pt cut in the offline selection.
%
%The extrapolation in the variable \alphat is tested through a
%dedicated set of closure tests, described in Sec.~\ref{sec:bkgd-syst},
%which demonstrate that the different \alphat requirements for the \mj
%and \mmj control samples and signal region have no significant
%systematic bias on the prediction. That the \alphat variable
%introduces no acceptance bias for processes with genuine \met is due
%to the accurate modelling by the CMS simulation of such processes,
%namely W + jets, \ttbar, and \znunu\ + jets. Background estimates for
%these processes are provided by the \mj and \mmj samples, as
%identified at the beginning of this Section. Importantly, the same
%cannot be said for QCD multijet events, as in this case the only
%events that survive the \alphat cut are pathological cases in which a
%jet is severely mismeasured or even lost due to detector
%inefficiencies. Such effects certainly do change the event kinematics
%and, in these cases, an \alphat cut will selectively choose events
%with particular kinematic features and topologies. For these
%pathological cases, one cannot rely on MC to model correctly the
%behaviour and therefore the \alphat acceptance. This is a crucial
%distinction to be made between QCD multijet events and processes with
%significant genuine \met. The assumption is that processes with
%genuine \met are selected by the \alphat variable based on the
%escaping invisible particle(s) rather than any pathological effects.

%\subsection{Distributions from the data control samples\label{sec:est-control-samples}}
%
%Distributions of key variables for the \mj, \mmj, and \gj samples are
%shown in below in Figs.~\ref{fig:mu-distr}, \ref{fig:mumu-distr},
%and~\ref{fig:phot-distr}, respectively. The first two figures show
%the (leading) muon \pt and isolation distributions, along with the
%leading jet \pt, \scalht, \mht and \alphat distributions. For the \gj
%sample, the photon \pt distribution is%and isolation distributions are
%shown in place of the corresponding muon distribution.%s.
%No requirement is made on the number of b-jets per event. In general,
%the agreement between data and simulation is good, giving confidence
%that the samples are well understood. The MC distributions highlight
%the composition of each sample. The contribution from QCD multijet
%events is expected to be negligible.
%
%Figure~\ref{fig:jet-bjet} shows the jet and b-jet multiplicity
%distributions for the \mj, \mmj, and \gj control samples, as defined
%in Sec.~\ref{sec:def-control-samples} and following the requirement
%$\scalht > 375\gev$. An accurate modelling of the jet multiplicity in
%data is achieved for all samples. The b-jet distributions demonstrate
%the changing background composition as a function of the number of
%b-jets. For the \mj and \mmj samples, the requirement of zero b-jets
%results in sub-samples that are rich in W and Z bosons, respectively,
%with little contamination from \ttbar. The \ttbar background becomes
%dominant in the \mj sample when exactly one b-jet is required. The
%requirement of up to two b-tags per event significantly suppresses all
%processes except for \ttbar production. Requiring at least three
%b-tags also suppresses \ttbar production. In the case of the \mmj
%sample, some contamination from \ttbar is observed for $\nb =
%1$. Requiring more than one b-jet significantly suppresses the yields
%in both the \mmj and \gj samples.
%
%\clearpage
%\begin{figure}[!h]
%  \centering
%  \subfigure[Muon \pt.]{
%    \includegraphics[width=0.4\textwidth]{figures/data-mc/v1/muon/Stack_muPt_Muon_all_OneMuon_-1To-2b_log}
%  } 
%  \subfigure[Muon isolation.]{
%    \includegraphics[width=0.4\textwidth]{figures/data-mc/v1/muon/Stack_muonIso_Muon_all_OneMuon_-1To-2b_log}
%  } \\
%  \subfigure[Transverse mass.]{
%    \includegraphics[width=0.4\textwidth]{figures/data-mc/v1/muon/Stack_PFMTmu_Muon_all_OneMuon_-1To-2b_log}
%  } 
%  \subfigure[\scalht.]{
%    \includegraphics[width=0.4\textwidth]{figures/data-mc/v1/muon/Stack_HT_Muon_all_OneMuon_-1To-2b_log}
%  } \\
%  \subfigure[\mht.]{
%    \includegraphics[width=0.4\textwidth]{figures/data-mc/v1/muon/Stack_MHT_Muon_all_OneMuon_-1To-2b_log}
%  } 
%  \subfigure[\alphat.]{
%    \includegraphics[width=0.4\textwidth]{figures/data-mc/v1/muon/Stack_AlphaT_Muon_all_OneMuon_-1To-2b_log}
%  } 
%  \caption{Data--MC comparisons of key variables for the \mj control
%    sample, for the region $\scalht > 275\GeV$. Bands represent the
%    uncertainties due to the limited size of MC samples. No
%    requirement on the number of b-jets per event is made.}
%  \label{fig:mu-distr}
%\end{figure}
%
%\begin{figure}[!h]
%  \centering
%  \subfigure[Leading muon \pt.]{
%    \includegraphics[width=0.4\textwidth]{figures/data-mc/v1/mumu/Stack_muPt_Muon_all_DiMuon_-1To-2b_log}
%  } 
%  \subfigure[Muon isolation.]{
%    \includegraphics[width=0.4\textwidth]{figures/data-mc/v1/mumu/Stack_muonIso_Muon_all_DiMuon_-1To-2b_log}
%  } \\
%  \subfigure[Di-muon invariant mass.]{
%    \includegraphics[width=0.4\textwidth]{figures/data-mc/v1/mumu/Stack_Zmass_Muon_all_DiMuon_-1To-2b_log}
%  } 
%  \subfigure[\scalht.]{
%    \includegraphics[width=0.4\textwidth]{figures/data-mc/v1/mumu/Stack_HT_Muon_all_DiMuon_-1To-2b_log}
%  } \\
%  \subfigure[\mht.]{
%    \includegraphics[width=0.4\textwidth]{figures/data-mc/v1/mumu/Stack_MHT_Muon_all_DiMuon_-1To-2b_log}
%  } 
%  \subfigure[\alphat.]{
%    \includegraphics[width=0.4\textwidth]{figures/data-mc/v1/mumu/Stack_AlphaT_Muon_all_DiMuon_-1To-2b_log}
%  } 
%  \caption{Data--MC comparisons of key variables for the \mmj control
%    sample, for the region $\scalht > 275\GeV$. Bands represent the
%    uncertainties due to the limited size of MC samples. No
%    requirement on the number of b-jets per event is made.}
%  \label{fig:mumu-distr}
%\end{figure}
%
%\begin{figure}[!h]
%  \centering
%  \subfigure[Photon \pt.]{
%    \includegraphics[width=0.4\textwidth]{figures/data-mc/v1/photon/Stacked_PhotonPt_all_Photon_375_upwards}
%  } 
%  \subfigure[\scalht.]{
%    \includegraphics[width=0.4\textwidth]{figures/data-mc/v1/photon/Stacked_HT_after_alphaT_55_all_Photon_375_upwards}
%  } \\
%  \subfigure[\mht.]{
%    \includegraphics[width=0.4\textwidth]{figures/data-mc/v1/photon/Stacked_MHT_after_alphaT_55_all_Photon_375_upwards}
%  } 
%  \subfigure[\alphat.]{
%    \includegraphics[width=0.4\textwidth]{figures/data-mc/v1/photon/Stacked_AlphaT_all_Photon_375_upwards}
%  } 
%  \caption{Data--MC comparisons of key variables for the \gj control
%    sample, for the region $\scalht > 375\GeV$. Bands represent the
%    uncertainties due to the limited size of MC samples. No
%    requirement on the number of b-jets per event is made.}
%  \label{fig:phot-distr}
%\end{figure}
%
%\begin{figure}[!h]
%  \centering
%  \subfigure[\njet for the \mj sample.]{
%    \includegraphics[width=0.4\textwidth]{figures/data-mc/v1/muon/Stack_ncommjet_Muon_all_OneMuon_-1To-2b_log}
%  } 
%  \subfigure[\njet for the \mmj sample.]{
%    \includegraphics[width=0.4\textwidth]{figures/data-mc/v1/mumu/Stack_ncommjet_Muon_all_DiMuon_-1To-2b_log}
%  } \\
%  \subfigure[\njet for the \gj sample.]{
%    \includegraphics[width=0.4\textwidth]{figures/data-mc/v1/photon/Stacked_JetMultiplicityAfterAlphaT_55_all_Photon_375_upwards}
%  } 
%  \subfigure[\nb for the \mj sample.]{
%    \includegraphics[width=0.4\textwidth]{figures/data-mc/v1/muon/Stack_nbjet_Muon_all_OneMuon_-1To-2b_log}
%  } \\
%  \subfigure[\nb for the \mmj sample.]{
%    \includegraphics[width=0.4\textwidth]{figures/data-mc/v1/mumu/Stack_nbjet_Muon_all_DiMuon_-1To-2b_log}
%  } 
%  \subfigure[\nb for the \gj sample.]{
%    \includegraphics[width=0.4\textwidth]{figures/data-mc/v1/photon/Stacked_Btag_Post_AlphaT_5_55_all_Photon_375_upwards}
%  } 
%  \caption{Data--MC comparison of the number of reconstructed jets
%    (top) and b-jets (bottom) per event in the (left) \mj sample,
%    (middle) \mmj sample, and (right) \gj control sample. Bands
%    represent the uncertainties due to the limited size of MC
%    samples.}\label{fig:jet-bjet}
%\end{figure}

%\subsection{Transfer factors and ``na\"ive'' predictions\label{sec:est-control-samples}}
%
%Appendix~\ref{app:tf} contains tables summarising the observed and
%expected yields from data and simulation, respectively, in the bins of
%the three control samples. Also listed are the expectations from
%simulation for the various background contributions in the signal
%region, along with the corresponding transfer factors. The yields
%are binned in \scalht, jet multiplicity and number of b-jets per
%event. The errors associated with the transfer factors reflect the
%uncertainty due to the finite size of the MC samples used to determine
%the factors. Any trigger inefficiency is also factored into the
%transfer factors (\ie, the trigger is effectively emulated and
%yields from the MC samples are corrected to account for any
%inefficiency). Also, all MC expectations are corrected to account for
%any discrepancies between data and MC for the efficiency and mistag
%rate of the b-tagging algorithm used, as described further in
%Sec.~\ref{sec:btag-eff-correction}. However, no systematic
%uncertainties on the transfer factors are quoted in the tables.
%
%The same tables also list ``na\"ive'' predicted yields obtained from
%each control sample for individual SM backgrounds in the signal region
%(\eg, W + jets and \ttbar from the \mj sample, and \znunu + jets from
%the \mmj and \gj samples). These predictions are given for
%illustrative purposes only. For the analysis result, the predictions
%for the total SM background are determined by a fit to the yields in
%the signal region and all three control samples, as described in
%Sec.~\ref{sec:statistics}. In addition to observed yields, the fit
%takes as input the transfer factors with their associated
%statistical and systematic uncertainties. 
%
%Illustrative predictions for the total SM background can be made for
%each bin in the signal region, by combining the individual
%predictions. One such combination can be made by using the individual
%predictions from the \mj and \mmj samples (or, alternatively, the
%predictions from the \mj and \gj samples), the result of which can be
%compared with the observed yields in the bins of the signal
%region. Predictions in this way are made for bins of the three
%exclusive b-tag categories requiring exactly zero and one b-tags per
%event. When requiring at least two b-tagged jets per event, only the
%\mj sample has sufficiently large yields to predict accurately the
%total SM background. The errors on the total SM predictions reflect
%statistical uncertainties only. It is again noted that these
%``na\"ive'' predictions are for illustration only, with the final SM
%expectations for all signal region bins given by the simultaneous fit
%to yields in all data samples. Table~\ref{tab:tf-summary} summarises
%the contents of the various sections in the Appendix that list tables
%containing observed yields, MC expectations, transfer factors, and
%``na\"ive'' predictions for individual and total SM backgrounds.
%
%\begin{table}[h!]
%  \caption{Each section in Appendix~\ref{app:tf} contains tables that
%    list: observed yields, MC expectations, transfer factors, and
%    ``na\"ive'' predictions for individual SM backgrounds, for each of
%    the \mj, \mmj, and \gj control samples; and total SM predictions
%    when combining the individual predictions from the \mj and \mmj
%    samples and the \mj and \gj samples, separately.}
%  \label{tab:tf-summary}
%  \centering
%  \footnotesize
%  \begin{tabular}{ llll }
%    \hline
%    Section            & \njet bin & \nb bin & Control samples used \\ [0.5ex]
%    \hline
%    \ref{app:23j0b}    & 2--3      & 0       & \mj, \mmj, \gj       \\
%    \ref{app:23j1b}    & 2--3      & 1       & \mj, \mmj, \gj       \\
%    \ref{app:23j2b1mu} & 2--3      & 2       & \mj                  \\
%    \ref{app:4j0b}     & $\geq$4   & 0       & \mj, \mmj, \gj       \\
%    \ref{app:4j1b}     & $\geq$4   & 1       & \mj, \mmj, \gj       \\
%    \ref{app:4j2b1mu}  & $\geq$4   & 2       & \mj                  \\
%    \ref{app:4j3b1mu}  & $\geq$4   & 3       & \mj                  \\
%    \ref{app:4j4b1mu}  & $\geq$4   & $\geq$4 & \mj                  \\
%    \hline
%  \end{tabular}
%\end{table}


%\begin{table}[h!]
%  \caption{}
%  \label{tab:}
%  \centering
%  \footnotesize
%  \begin{tabular}{ lll }
%    \hline
%    \hline
%    \njet bin & \nb bin & Control samples used \\ [0.5ex]
%    \hline
%    2--3      & 0       & \mj, \mmj, \gj       \\
%    2--3      & 1       & \mj, \mmj, \gj       \\
%    2--3      & 2       & \mj                  \\
%    $\geq$4   & 0       & \mj, \mmj, \gj       \\
%    $\geq$4   & 1       & \mj, \mmj, \gj       \\
%    $\geq$4   & 2       & \mj                  \\
%    $\geq$4   & 3       & \mj                  \\
%    $\geq$4   & $\geq$4 & \mj                  \\
%    \hline
%    \hline
%  \end{tabular}
%\end{table}

